% Generated draft: requires mathematical and source review.
\documentclass[11pt]{article}
\usepackage[margin=25mm]{geometry}
\usepackage{fontspec}
\setmainfont{Latin Modern Roman}
\usepackage{amsmath,amssymb,mathtools,unicode-math}
\setmathfont{Latin Modern Math}
\usepackage{enumitem,needspace,xurl,hyperref}
\setlength{\parindent}{0pt}
\setlength{\parskip}{5pt}
\setlength{\emergencystretch}{3em}
\newcommand{\N}{\mathbb{N}}
\newcommand{\Z}{\mathbb{Z}}
\newcommand{\Q}{\mathbb{Q}}
\newcommand{\R}{\mathbb{R}}
\newcommand{\C}{\mathbb{C}}
\newcommand{\F}{\mathbb{F}}
\newcommand{\A}{\mathbb{A}}
\newcommand{\PP}{\mathbb{P}}
\newcommand{\E}{\mathbb{E}}
\newcommand{\eps}{\varepsilon}
\newcommand{\dd}{\,\mathrm{d}}
\begin{document}
\section*{UnsolvedMath ID: 30000185}
\textbf{Source problem code:} catalogue entry
\subsection{Definitions and mathematical statement}
The source asks whether the number \(m(X,Y)\) of rational maps from a curve \(X\) to a curve \(Y\) is polynomially bounded in the genus \(g(X)\). It gives no field, smoothness or projectivity hypotheses, genus restriction on \(Y\), or convention for which maps are counted. These omissions affect finiteness before they affect a quantitative bound.

Under the usual illustrative convention of smooth projective connected curves over \(\C\), the genus is \(g(X)=\dim_{\C}H^0(X,\Omega_X^1)\). A rational map is a morphism on a nonempty Zariski-open subset, with two representatives identified when they agree on a nonempty open subset. Let \(\mathcal R(X,Y)\) stand for the collection intended by the source, and write \(m(X,Y)=|\mathcal R(X,Y)|\). This notation leaves unresolved whether \(\mathcal R\) means all rational maps, only nonconstant maps, or maps modulo a specified equivalence; the draft does not silently substitute one of these choices.

For a specified admissible class \(\mathcal A\) of pairs \((X,Y)\), a uniform polynomial-bound interpretation is
\[
 \exists C>0\ \exists d\in\Z_{\geq0}\ \forall (X,Y)\in\mathcal A,
 \qquad m(X,Y)\leq C(1+g(X))^d.
\]
If the target is fixed instead, the quantifiers might be \(\forall Y\ \exists C_Y,d_Y\ \forall X\), or the degree might be required to be absolute while only the constant depends on \(Y\). The source does not distinguish these requests. The term \(1+g(X)\) is only a convenient convention for expressing a polynomial growth bound at all nonnegative genera.

If all constant maps between complex curves are counted, the points of \(Y(\C)\) already supply infinitely many maps. Merely excluding constants still leaves an essential target-genus issue: unrestricted maps to \(\PP^1\) and elliptic targets need not form finite collections. These are missing-hypothesis warnings, not a replacement conjecture or a claimed resolution of the intended problem. Restrictions such as a hyperbolic target, a degree bound or an automorphism quotient cannot be inferred from the extracted wording.

The literature assessment checked 2026-08-21 says only a lead-in survives and the conjecture itself is missing, whereas the supplied original and clean fields both contain the displayed map-counting question. The verification note reports restoration of notation only. This source mismatch and the supplied open label require review; neither furnishes the missing hypotheses or a current theorem-status conclusion.
\subsection{Short English statement}
The intended question asks for a polynomial bound, in the source genus, on a count of maps to another curve. A meaningful finite count needs the field, curve hypotheses, nonconstant-map convention, target restrictions and uniformity of the bound. These data are absent, and the literature note appears to describe a different or earlier extraction.
\subsection{Sources}
\begin{itemize}
\item Local UnsolvedMath registry, permanent ID 30000185, source code catalogue entry, ``Polynomial Bounds for Rational Maps Between Curves,'' including the original and cleaned question, notation-restoration verification and conflicting extraction assessment.
\item \emph{Komplexe Algebraische Geometrie (2005)}, Oberwolfach Reports, pp.\ 407--463, DOI 10.4171/owr/2005/07: \url{https://doi.org/10.4171/owr/2005/07}. This is the sole supplied literature reference, labelled ``original OWR report.''
\item Literature triage checked 2026-08-21; registry record updated 2026-09-16. No external source has been visited or independently reviewed here.
\end{itemize}
\end{document}
